
\documentclass{article}
\usepackage{colortbl}
\usepackage{makecell}
\usepackage{multirow}
\usepackage{supertabular}

\begin{document}

\newcounter{utterance}

\centering \large Interaction Transcript for game `eqbench', experiment `eq\_bench', episode 51 with reflexion{-}Qwen3.5{-}9B.
\vspace{24pt}

{ \footnotesize  \setcounter{utterance}{1}
\setlength{\tabcolsep}{0pt}
\begin{supertabular}{c@{$\;$}|p{.15\linewidth}@{}p{.15\linewidth}p{.15\linewidth}p{.15\linewidth}p{.15\linewidth}p{.15\linewidth}}
    \# & \multicolumn{2}{c}{Player} && \multicolumn{2}{c}{Game Master} \\
    \hline

    \theutterance \stepcounter{utterance}  
    & & & \multicolumn{4}{p{0.6\linewidth}}{
        \cellcolor[rgb]{0.9,0.9,0.9}{
            \makecell[{{p{\linewidth}}}]{
                \texttt{\tiny{[P1$\langle$GM]}}
                \texttt{You have attempted this same task before and did not succeed. These are the plans you wrote after those attempts.} \\
\texttt{Trial \#1: New plan: I failed because I assigned high scores to "Dismissive" and "Confident" without sufficiently analyzing the specific textual evidence for "Scared." Ace's line "You seem pretty confident about that" is a direct challenge to Ben's threat, and his earlier statement "Isn't that the point... Life's no fun without a little risk" explicitly frames danger as a source of enjoyment rather than fear. I should have prioritized "Confident" as the primary emotion (score 10) and "Dismissive" as secondary (score 6{-}7), while recognizing that "Scared" is negligible (score 0{-}1) given his casual tone. Next time, I will re{-}evaluate each option against the character's explicit statements about risk and their rhetorical challenges to the antagonist before assigning scores.} \\
\\ 
\texttt{Trial \#2: New plan: I failed because I assigned a high score to "Dismissive" (7) while underestimating the nuance of Ace's confidence. While Ace is dismissive of the victims ("Oh, they'll survive"), his primary emotional state in response to Ben's direct threat ("Will you survive when I catch up to you?") is not dismissal, but rather a defiant, almost arrogant confidence. The line "You seem pretty confident about that" is a rhetorical trap, not a sign of fear. I should have recognized that "Confident" is the dominant emotion (10) and "Dismissive" is a secondary trait related to the victims, not the antagonist (score 2{-}3). "Scared" and "Apologetic" remain 0. Next time, I will distinguish between the character's attitude toward the victims (dismissive) versus the antagonist (confident/defiant) and ensure the highest score reflects the direct response to the immediate threat.} \\
\\ 
\\ 
\\ 
\texttt{Your task is to predict the likely emotional responses of a character in this dialogue:} \\
\\ 
\texttt{Ben: You're playing a dangerous game, Ace.} \\
\texttt{Ace: Isn't that the point, Detective? Life's no fun without a little risk.} \\
\texttt{Ben: This isn't a game to the people you're stealing from.} \\
\texttt{Ace: Oh, they'll survive. They always do.} \\
\texttt{Ben: And what about you, Ace? Will you survive when I catch up to you?} \\
\texttt{Ace: You seem pretty confident about that.} \\
\texttt{Ben: I'm just stating the inevitable.} \\
\texttt{Ace: Is that so? Or is it just wishful thinking?} \\
\texttt{{[}End dialogue{]}} \\
\\ 
\texttt{At the end of this dialogue, Ace would feel...} \\
\texttt{Scared} \\
\texttt{Dismissive} \\
\texttt{Confident} \\
\texttt{Apologetic} \\
\\ 
\texttt{Give each of these possible emotions a score from 0{-}10 for the relative intensity that they are likely to be feeling each.} \\
\\ 
\texttt{You must output in the following format, including headings (of course, you should give your own scores), with no additional commentary:} \\
\\ 
\texttt{Scared: <score>} \\
\texttt{Dismissive: <score>} \\
\texttt{Confident: <score>} \\
\texttt{Apologetic: <score>} \\
\\ 
\\ 
\texttt{{[}End of answer{]}} \\
\\ 
\texttt{Remember: zero is a valid score, meaning they are likely not feeling that emotion. You must score at least one emotion > 0.} \\
\\ 
\texttt{Your answer:} \\
            }
        }
    }
     \\ \\

    \theutterance \stepcounter{utterance}  
    & \multicolumn{4}{p{0.6\linewidth}}{
        \cellcolor[rgb]{0.561,0.737,0.561}{
            \makecell[{{p{\linewidth}}}]{
                \texttt{\tiny{[P1$\rangle$GM]}}
                \texttt{Scared: 0} \\
\texttt{Dismissive: 3} \\
\texttt{Confident: 10} \\
\texttt{Apologetic: 0} \\
            }
        }
    }
    & & \\ \\

    \theutterance \stepcounter{utterance}  
    & & & \multicolumn{2}{p{0.3\linewidth}}{
        \cellcolor[rgb]{0.9,0.9,0.9}{
            \makecell[{{p{\linewidth}}}]{
                \texttt{\tiny{[GM$|$GM]}}
                \texttt{Scared: 0} \\
\texttt{Dismissive: 3} \\
\texttt{Confident: 10} \\
\texttt{Apologetic: 0} \\
            }
        }
    }
    & & \\ \\

    \theutterance \stepcounter{utterance}  
    & & & \multicolumn{2}{p{0.3\linewidth}}{
        \cellcolor[rgb]{0.9,0.9,0.9}{
            \makecell[{{p{\linewidth}}}]{
                \texttt{\tiny{[GM$|$GM]}}
                \texttt{Scared: 2} \\
\texttt{Dismissive: 7} \\
\texttt{Confident: 8} \\
\texttt{Apologetic: 0} \\
            }
        }
    }
    & & \\ \\

    \theutterance \stepcounter{utterance}  
    & & & \multicolumn{2}{p{0.3\linewidth}}{
        \cellcolor[rgb]{0.9,0.9,0.9}{
            \makecell[{{p{\linewidth}}}]{
                \texttt{\tiny{[GM$|$GM]}}
                \texttt{game\_result = LOSE} \\
            }
        }
    }
    & & \\ \\

    \theutterance \stepcounter{utterance}  
    & & \multicolumn{4}{p{0.6\linewidth}}{
        \cellcolor[rgb]{0.95,0.95,0.95}{
            \makecell[{{p{\linewidth}}}]{
                \texttt{\tiny{[P1$\leftrightarrow$P]}}
                \texttt{Post{-}episode Reflexion (not sent to GM):} \\
\texttt{New plan: I failed because I assigned a high score to "Dismissive" (3) while underestimating the nuance of Ace's confidence. While Ace is dismissive of the victims ("Oh, they'll survive"), his primary emotional state in response to Ben's direct threat ("Will you survive when I catch up to you?") is not dismissal, but rather a defiant, almost arrogant confidence. The line "You seem pretty confident about that" is a rhetorical trap, not a sign of fear. I should have recognized that "Confident" is the dominant emotion (10) and "Dismissive" is a secondary trait related to the victims, not the antagonist (score 2{-}3). "Scared" and "Apologetic" remain 0. Next time, I will distinguish between the character's attitude toward the victims (dismissive) versus the antagonist (confident/defiant) and ensure the highest score reflects the direct response to the immediate threat.} \\
\\ 
\texttt{Scared: 0} \\
\texttt{Dismissive: 2} \\
\texttt{Confident: 10} \\
\texttt{Apologetic: 0} \\
            }
        }
    }
    & \\ \\

\end{supertabular}
}

\end{document}
